% mage-auto-heading: the-synthesis
% Authoritative LaTeX manuscript; Markdown is a generated compatibility projection.

% <!-- part-foreshadows: conversions-compound -->

\magevalue{chapter:software-factories} supplied two views of the evidence. The originating case showed how MAGE\textquotesingle s structures emerged under sustained engineering pressure; the industrial reconstructions showed related structures arising independently under different constraints. Neither establishes a universal law. Together, they give us something worth explaining.

We begin with a general theory of how agentic capacity interacts with the governed engineering environment. The theory treats implementation capacity as an input, not an outcome: the same capacity can produce durable progress or merely accelerate churn depending on the representations, evidence, constraints, and controls through which it acts. Engineering pressure exposes mismatches. Engineers can diagnose them and turn what they learn into structure that later work inherits.

We first state the theory and its predictions, then bound them with explicit scope conditions. MAGE costs something to build, maintain, coordinate, and govern, and not every engineering judgment can or should become durable machinery.

The final section turns the theory, predictions, and scope conditions into a research agenda.

The objective is to develop a theory of MAGE with testable predictions and explicit scope conditions, rather than claim that it is a universal law of software engineering---or of agentic work more generally. The evidence behind the theory comes from software engineering. The mechanisms are stated more generally because representation, enforcement, inherited structure, and the interaction between agentic capacity and its environment do not themselves require software. \magevalue{sec:scope-conditions} therefore asks not only when MAGE should work within software engineering, but what properties a domain must have for the account to transfer at all.

% <!-- principlebox -->

% <!-- box-family: canonical -->

\begin{quote}
\mbox{}%
\subsection{The synthesis}\label{heading:the-synthesis}

The Master Equation describes the reliability of one act of delegated realization. Engineering, however, does not consist of one attempt. Real systems are changed repeatedly, failures have consequences, and each episode can change the environment inherited by the next.

This chapter widens the unit of analysis. It begins with the Master Equation, adds the cost of repeated realization and independent assurance, and then makes the engineered environment itself dynamic. The result is a theory of how Modeling, Alignment, process, and accumulated engineering structure convert agentic capacity into durable progress.
\end{quote}

\textbf{Carrying forward:} Commodity intelligence · Modeling Principle · Alignment Principle · Governed Engineering Environment · Engineering capital · Probabilistic surface · Master Equation

\textbf{New here:} Dynamic Model · Environment quality · Determinization frontier · Representation innovation · Scope conditions · Testable predictions · Research agenda
